\sectionCenteredToc{ЗАКЛЮЧЕНИЕ}

В ходе выполнения дипломного проекта успешно достигнута поставленная цель — разработан, программно реализован и исследован гибридный асимметричный программный комплекс для адаптивного сжатия и восстановления цифровых изображений на базе каскадного дискретного вейвлет-преобразования (DWT) и легковесных нейросетевых моделей сверхразрешения (ESPCN, FSRCNN).

В процессе работы были решены все поставленные задачи:
\begin{enumerate}
    \item Проведен сравнительный анализ существующих методов компрессии растровой графики, выявивший ограничения алгоритмов на базе DCT (JPEG) при экстремальных степенях сжатия и обосновавший необходимость перехода к гибридным кодекам.
    \item Сформулировано и обосновано техническое задание на разработку комплекса с определением функциональных требований, критериев надежности и совместимости в строгом соответствии со стандартами ГОСТ~28195-89 и СТБ~ИСО/МЭК~12207-2003.
    \item Разработана структурная схема системы, реализующая асимметричный подход Edge-to-Cloud к распределению нагрузки: энергоэффективное математическое кодирование на стороне периферийного узла-отправителя и ресурсоемкий тензорный инференс на стороне облачного сервера.
    \item Спроектированы и реализованы инновационные алгоритмы конвейера обработки: модуль адаптивной маршрутизации файлов, механизм динамической экстракции гибридных метаданных (цветовых паспортов и карт контуров Кенни), а также тайловая маршрутизация памяти.
    \item Проведены экспериментальные исследования разработанного веб-приложения, доказавшие эффективность алгоритма: при сжатии трафика на 80–95\,\% декодер успешно реконструирует оригинальную структуру со стабильно высокими показателями качества (PSNR выше 30~дБ, SSIM выше 0,80).
\end{enumerate}

Разработанный программный комплекс полностью функционален, обладает отказоустойчивостью и готов к практическому внедрению. Предложенный асимметричный подход имеет высокую практическую значимость для оптимизации мультимедийного трафика в сетях с ограниченной пропускной способностью, гарантируя радикальное снижение битрейта при сохранении визуальной полноты передаваемых данных.